\section{Descent of the difficulty and klt termination}
\label{sec:difficulty}

Return to the prepared terminal chain of
Section~\ref{sec:difficulty-setup}.  Its labels $S_h$ persist, their
coefficients are nonincreasing, and their default echo weights $d_h$
are constant.  The integer $\mathcal D$ is the finite corrected sum
in Definition~\ref{def:difficulty}.

The branch theorem now supplies the lower bound isolated in
\eqref{eq:difficulty-local-deficit}.  To obtain termination, we must
also compare this integer across every move of the chain.  The
normalization-rank correction was chosen so that its change cancels
exactly the change in the subtracted branch defaults.  We prove this
identity first, then show that each positive-side surface downstairs
forces an integer weight to drop between consecutive junctions.

\subsection{Divisor ranks on the normalizations}

We record explicitly the homological fact which will cancel the
changes of the defaults under a small flip.

\begin{lemma}
\label{difficulty:normalization-independence}
Let \(q:S\to T\) be a projective birational morphism of normal
projective threefolds.  The cycle classes in \(H_4(S,\mathbb Q)\)
of the prime divisors contracted by \(q\) are linearly independent.
\end{lemma}

\begin{proof}
Take a projective resolution \(r:\widetilde S\to S\) which is an
isomorphism over \(U=S_{\mathrm{reg}}\), and put
\(E=r^{-1}(S\setminus U)\).  Normality gives
\(\dim(S\setminus U)\leq1\), so \(\dim E\leq2\).
Suppose a rational combination of the contracted divisors has zero
class in \(H_4(S,\mathbb Q)\).  The corresponding combination of
their strict transforms has zero restriction to
\(H_4^{\mathrm{BM}}(U,\mathbb Q)\).  The localization sequence
\[
 H_4(E,\mathbb Q)\longrightarrow
 H_4(\widetilde S,\mathbb Q)\longrightarrow
 H_4^{\mathrm{BM}}(U,\mathbb Q)
\]
is exact.  Its first group is generated by the fundamental classes
of the irreducible surface components of \(E\); components of
dimension at most one contribute nothing in degree four.  Consequently
one can correct the strict-transform combination by a rational
combination of resolution-exceptional divisors to obtain a divisor
with zero homology class on \(\widetilde S\).

Every divisor in this corrected combination is exceptional over
\(T\).  This is true for the original strict transforms by
assumption.  For a resolution-exceptional divisor it follows because
its image in \(S\) has dimension at most one, and hence so does
its image in \(T\).  On the smooth projective threefold
\(\widetilde S\), a homologically trivial divisor is numerically
trivial.  Applying the negativity lemma to this exceptional divisor
and to its negative shows that it is zero as a divisor.  None of
the strict transforms is resolution-exceptional, so all original
coefficients vanish.
\end{proof}

Consider now one of the small terminal flips of our chain, written
\(Y^-\to Z\leftarrow Y^+\), with the full bad sets
\(W^-,W^+\).  For a label \(S_h\), let \(T_h\) be the
normalization of its common image in \(Z\).  Since the ambient maps
are small, the two maps
\[
 (S_h^-)^\nu\longrightarrow T_h
 \longleftarrow(S_h^+)^\nu
\]
are projective birational morphisms.  Let \(e_h^\pm\) be the number
of divisors in these normalizations lying over the corresponding
bad set.  They are exactly the divisors counted by
\(r_{V,h}\) for surfaces \(V\subset W^\pm\).

Indeed, normalization is finite and so preserves the dimension of
these images in \(S_h^\pm\).  A divisor in the normalization has
a surface image.  The bad locus on the ambient base has dimension
at most one by Lemma~\ref{lem:graph}; consequently all these divisors
are contracted over \(T_h\).  Conversely, outside the bad set
the ambient maps, and hence the normalizations of their labeled
divisors, identify.  Thus all exceptional divisors of the displayed
morphisms occur in this list.

Restriction to the common open identifies the spans of the
restricted divisor cycle classes: every divisor on that open has
a closure on either projective normalization.  By localization,
the kernels of these restrictions are the spans of the bad-locus
divisors.  Lemma~\ref{difficulty:normalization-independence} computes
their dimensions as \(e_h^\pm\).  We obtain the exact equality
\begin{equation}
\label{difficulty:rank-difference}
 \rho_*((S_h^-)^\nu)-\rho_*((S_h^+)^\nu)
   = e_h^- - e_h^+.
\end{equation}
This equality concerns divisor cycle classes on the normalizations;
it does not require a Picard-number identity for the ambient flip.

\subsection{Boundedness and monotonicity}

\begin{proposition}
\label{prop:difficulty}
On the prepared terminal chain, the integer
\(\mathcal D(Y,\Theta)\) has the following properties.
\begin{enumerate}[label=\textup{(\roman*)}]
\item It is nonnegative at every crepant junction.
\item It is nonincreasing under every coefficient decrease and
every small terminal flip.
\item If the weight of one exceptional valuation strictly decreases
in one of these moves, then the difficulty strictly decreases.
\end{enumerate}
No separate nonnegativity assertion is required for the intermediate
models in a bridge.
\end{proposition}

\begin{proof}
At a crepant junction \(p:Y\to X_i\), let \(Q\) be the reduced
union of the varying labels.  If this union is empty,
Lemma~\ref{lem:ceiling-error} already makes every surface bracket
nonnegative, and all other terms of \eqref{eq:difficulty} are
nonnegative.  Otherwise each component of \(Q\) has positive
coefficient and, by Proposition~\ref{prop:codim-two},
\(\dim p(Q)\leq1\).  For every surface in \(Q\), its branch
number in Theorem~\ref{thm:branch} equals
\(r_V^{\mathrm{var}}\).  Surfaces outside \(Q\) have no negative
allowance in \eqref{eq:local-ceiling-bound}.

For a varying label \(b_h>0\), so \(d_h\geq w(b_h)\geq1\).
Moreover, every valuation in the right-hand count of
Theorem~\ref{thm:branch} has a codimension-three center and weight
at least one, and thus is counted in the first term of
\eqref{eq:difficulty}.  It follows from that theorem that
\[
 \sum_{V\subset Q}(r_V^{\mathrm{var}}-1)
 \leq
 \sum_{h\text{ varying}}d_h\rho_*(S_h^\nu)
 +\sum_{\operatorname{codim}_Y c_Y(v)\geq3}w_v.
\]
Substitution in \eqref{eq:difficulty-local-deficit} proves (i).
The valuations on the right have codimension-three centers, so none
has been reused as a valuation in a surface bracket.

For a coefficient decrease the variety, all centers, all branch
numbers, and all normalization ranks remain fixed.  The boundary
removed is effective and \(\mathbb Q\)-Cartier; hence every
discrepancy is nondecreasing and every weight is nonincreasing.
The \(d_h\) remain unchanged by our preparation.  To check that
this comparison is a finite one, take the union of the exceptional
surface centers in Lemma~\ref{lem:echoes} for the two boundaries.
Outside this finite set a nonzero local contribution consists of
the echoes on a single label.  Each of their finitely many possibly
positive weights is nonincreasing, while their sum is the unchanged
\(d_h\).  Thus each such echo weight is unchanged.  The
codimension-at-least-three contributions, and the contributions at
the finitely many exceptional surface centers, can therefore be
compared as finite sums with unchanged corrections.  Their weights
only decrease; a strict decrease of any exceptional valuation
strictly decreases the total.  This proves (ii) and (iii) for a
coefficient move.

For a small flip, all prime labels and all exceptional divisorial
valuations correspond on the two sides.  By
Lemma~\ref{lem:graph}, the set
\[
\begin{split}
 \mathcal A
 &=\{v:\ v\text{ is exceptional and }c_{Y^-}(v)\subset W^-\}\\
 &=\{v:\ v\text{ is exceptional and }c_{Y^+}(v)\subset W^+\}
\end{split}
\]
is the same on both sides.  Set
\[
 \mathcal A_0=
 \{v\in\mathcal A:\ w_v(Y^-,\Theta^-)+w_v(Y^+,\Theta^+)>0\}.
\]
This is finite.  Each bad set has dimension at most two and has
finitely many surface components; a valuation with a surface center
contained in it has one of these centers.  Lemma~\ref{lem:echoes}
gives finiteness at each of them and finiteness of all positive-weight
terms with centers of codimension at least three.

The generic centers not contained in the bad sets identify on the
common open.  Their discrepancies, local corrections, and labels
therefore agree.  Write \(w_v^\pm=w_v(Y^\pm,\Theta^\pm)\).
Separating the finite bad-set contributions in
\eqref{eq:difficulty} gives
\begin{align}
\label{difficulty:flip-difference}
 \mathcal D(Y^-,\Theta^-)-\mathcal D(Y^+,\Theta^+)
 ={}&\sum_{v\in\mathcal A_0}(w_v^- - w_v^+)
       -\sum_h d_h(e_h^- - e_h^+) \\
    &+\sum_h d_h
      \bigl(\rho_*((S_h^-)^\nu)-\rho_*((S_h^+)^\nu)\bigr)
 \notag\\
 ={}&\sum_{v\in\mathcal A_0}(w_v^- - w_v^+).
 \notag
\end{align}
The second equality is exactly
\eqref{difficulty:rank-difference}.  Discrepancy monotonicity makes
every summand in the last expression nonnegative.  A strict weight
drop gives a positive summand and hence a strict drop of the
difficulty.  Valuations whose centers change from a surface to a
curve, or conversely, are still counted once in the raw part on each
side; they merely move between the first two displayed sums of
\eqref{eq:difficulty}.  Positive weights becoming zero are included
in \(\mathcal A_0\).  Thus neither kind of change is lost in the
comparison.  This completes (ii) and (iii).
\end{proof}

\begin{remark}
The local bracket at a bad surface may cancel the weight of a
generic echo.  Formula~\eqref{difficulty:flip-difference} shows why
this cannot absorb a strict decrease: the change of that default is
matched exactly by the normalization-rank term.  This remains true
when a varying coefficient tends to zero while its positive default
stays constant.  An intermediate difficulty may initially have been
defined as a signed integer; Proposition~\ref{prop:difficulty}
then bounds it below by the nonnegative value at the next junction.
\end{remark}

\subsection{Detecting positive-side surfaces}

\begin{theorem}
\label{thm:klt-termination}
Every program in Definition~\ref{def:program} starting with a
projective klt fourfold pair over \(\mathbb C\) with rational boundary has finitely
many birational steps.  No pseudo-effectivity hypothesis is needed.
\end{theorem}

\begin{proof}
Suppose that such a program has infinitely many birational steps.
Lemma~\ref{lem:cycle-rank}, applied in cycle dimension three,
removes all divisor-losing steps after a finite prefix.  The
remaining program is small.  Construct its terminal chain by
Proposition~\ref{prop:terminal-chain}, and make the truncations in
Proposition~\ref{prop:codim-two} and
Lemma~\ref{difficulty:preparation}.  By
Proposition~\ref{prop:difficulty}, the difficulties at its crepant
junctions are nonnegative integers and are nonincreasing.

Consider a downstairs step with an irreducible surface
\(V\subset W_i^+\).  Proposition~\ref{prop:codim-two} says that
the ambient variety is smooth at \(\eta_V\) and that its
minimal log discrepancy there exceeds one.  The ordinary generic
blowup of \(V\) defines a normalized valuation \(v\) whose
positive-side discrepancy is
\[
 a^+=a(v;X_{i+1},B_{i+1})
     =2-\sum_hm'_h b'_h>1,
\]
where \(m'_h\) are boundary multiplicities at \(\eta_V\) and
\(b'_h\) are the original boundary coefficients carried by strict
transform downstairs.  Effectiveness gives \(a^+\leq2\), and
the choice of \(N\) gives
\[
 M:=N(2-a^+)=N\sum_hm'_h b'_h\in\mathbb Z_{\geq0}.
\]
By the full bad-set strictness in Lemma~\ref{lem:graph},
\(a^-:=a(v;X_i,B_i)<a^+\).  Consequently
\begin{equation}
\label{eq:difficulty-unit-drop}
 w(2-a^-)\geq M+1,
 \qquad w(2-a^+)=M.
\end{equation}
In particular the argument includes \(a^+=2\), when the
positive-side weight is zero.  Strict discrepancy increase alone
would not imply this unit drop without the integral final deficit.

The valuation \(v\) is exceptional downstairs on both sides,
since the step is small.  It does not belong to the stabilized list
\(I_0\): its positive-side discrepancy exceeds one, whereas
\(I_0\) is the common list of exceptional valuations of discrepancy
at most one.  The terminal junctions extract exactly \(I_0\), so
\(v\) is exceptional on both junction models.  Every intervening
birational move is small, and a coefficient change does not alter
the model.  Hence \(v\) stays exceptional throughout this finite
bridge.  Crepancy at its endpoints identifies its junction weights
with those in \eqref{eq:difficulty-unit-drop}.  At least one move
in the bridge therefore strictly lowers its weight.  By
Proposition~\ref{prop:difficulty}, the difficulty strictly decreases
between these consecutive junctions.

A nonincreasing sequence of nonnegative integers has only finitely
many strict decreases.  Thus only finitely many downstairs steps
have positive-side bad surfaces.  Discard them.  On the remaining
small tail, Lemma~\ref{lem:cycle-rank} in cycle dimension two
removes negative-side bad surfaces after finitely many further
steps.  Both bad loci of every subsequent step would then have
dimension at most one, contradicting the inequality
\(\dim W_i^-+\dim W_i^+\geq3\) of Lemma~\ref{lem:graph}.
This contradiction proves termination.
\end{proof}
